\section{Accessible categories and the relative classifying topos (Part I)}\label{sec:categoryI}

\subsection{The indexing convention and the actual equivalence}
We fix a universe and use ``small'' relative to it.  A category is
$\kappa$-filtered, for an infinite regular cardinal $\kappa$, when every
diagram having fewer than $\kappa$ arrows has a cocone.  An object $c$ is
$\kappa$-presentable when $\operatorname{Hom}(c,-)$ preserves small
$\kappa$-filtered colimits.  The manuscript's ordered sets ``grands devant
$\pi$'' have upper bounds for subsets of cardinal at most $\pi$; their
modern index is therefore $\kappa=\pi^+$, not $\pi$.
For a small category $C$, $\operatorname{Ind}_\kappa(C)$ is the full
subcategory of $[C^{\mathrm{op}},\mathbf{Set}]$ consisting of small
$\kappa$-filtered colimits of representables.
A full subcategory $C\subset E$ is dense when the restricted Yoneda
functor $N_C:E\to[C^{\mathrm{op}},\mathbf{Set}]$ is fully faithful.

\begin{theorem}[Proved: the corrected equivalence]\label{cat:I1}
For a locally small category $E$ with small $\kappa$-filtered colimits,
the following conditions are equivalent.
\begin{enumerate}
\item There is a small full subcategory of $\kappa$-presentable objects
such that every object of $E$ is a $\kappa$-filtered colimit of its objects.
\item $E\simeq\operatorname{Ind}_\kappa(C)$ for a small category $C$.
\item The $\kappa$-presentable objects form an essentially small dense
subcategory $E_\kappa$, and $E_\kappa/X$ is $\kappa$-filtered for every
$X\in E$.
\end{enumerate}
In $\operatorname{Ind}_\kappa(C)$ the $\kappa$-presentable objects are
exactly the retracts of representables.  If $E$ is cocomplete, condition
(3) may omit the last, slice-filteredness condition.  In that case it
also suffices that $E$ have a small dense subcategory of
$\kappa$-presentable objects.
\end{theorem}
\begin{proof}
Representables are $\kappa$-presentable because evaluation commutes with
colimits of presheaves.  If $X=\operatorname{colim}_I yc_i$ is itself
$\kappa$-presentable, its identity factors through some $yc_i$, so it is
a retract.  Conversely a retract of a presentable object is presentable:
the associated Hom functor is a retract of a functor preserving the
colimits in question.  Such retracts do belong to
$\operatorname{Ind}_\kappa(C)$.  For an idempotent $e$ on $yc$, take the
constant-object diagram on the ordinal $\kappa$ whose maps between
strictly increasing indices are $e$; its colimit is the image of $e$.
The ordinal is $\kappa$-filtered by regularity.  There is only a set of
idempotents on representables, which proves essential smallness.

For any presheaf $X$ which is a $\kappa$-filtered colimit of
representables, $C/X$ is $\kappa$-filtered.  Indeed, every family of
fewer than $\kappa$ elements, and all relations in a diagram between
them, is represented at a common stage, and the fewer than $\kappa$
required equalities hold at a further common stage.  That stage defines
a cocone in $C/X$.  The usual presheaf presentation
$X=\operatorname{colim}_{C/X}yc$ therefore gives the desired canonical
filtered presentation.  The same element argument shows that a
$\kappa$-filtered colimit of such presheaves again has a
$\kappa$-filtered category of elements.  Thus these colimits remain in
$\operatorname{Ind}_\kappa(C)$ and are computed pointwise.

Suppose (1) holds, with the chosen small category $C$.  For a
presentation $X=\operatorname{colim}_I c_i$, presentability gives
$N_CX=\operatorname{colim}_I yc_i$.  Hence, for $Y\in E$,
\[
 \operatorname{Hom}(N_CX,N_CY)
 =\lim_i\operatorname{Hom}(c_i,Y)
 =\operatorname{Hom}_E(X,Y).
\]
The comparison is induced by $N_C$, so this is full faithfulness.
Every filtered colimit of representables is the nerve of the
corresponding colimit in $E$.  We obtain (2).  The preceding paragraphs
then prove (3), replacing $C$ by its essentially small retract closure.
Conversely, under (3), density identifies $X$ with its canonical colimit
of $E_\kappa/X$; this colimit exists by the assumed filteredness.
This is (1).

Finally let $E$ be cocomplete.  Colimits of fewer than $\kappa$
presentable objects are presentable: their Hom functors are limits of
fewer than $\kappa$ set-valued functors, and $\kappa$-filtered colimits
of sets commute with such limits.  The latter assertion follows by
choosing a common stage for fewer than $\kappa$ elements and a further
stage for fewer than $\kappa$ equalities.  Thus $E_\kappa/X$ is
$\kappa$-filtered whenever $E_\kappa$ is essentially small.
Starting only with a small dense $C$ of presentables, close it under
colimits of diagrams with fewer than $\kappa$ arrows.  Iterating for
$\kappa$ stages gives a small category $C'$ of presentables, closed
under these colimits.  It is still dense: a natural map on $C'$ is
uniquely determined on $C$, and the maps from $C$ jointly detect its
components.  The category $C'/X$ is filtered as above.  Applying the
already proved argument to $C'$ gives (1) and the required smallness
of all presentables.
\end{proof}

This is exactly the extra condition written on 161-2/060, conditions
(i)--(iii) of the converse.  The parenthetical uncertainty on scan 055
must not erase it.  The proof also spells out why the cocomplete case
on scans 061--062 is different.

\subsection{What the cardinal estimate actually says}
\begin{lemma}[Proved: small directed subdiagrams]\label{cat:I2}
Let $I$ be a $\kappa$-directed poset, let $\mu\geq\kappa$ be infinite,
and assume $\mu^{<\kappa}=\mu$.  Every subset of $I$ of size at most
$\mu$ is contained in a $\kappa$-directed subset of size at most $\mu$.
Let $J_\mu$ be the poset of these directed subsets, ordered by inclusion.
It is $\mu^+$-directed, its union is $I$, and
\[
 \operatorname{cf}(J_\mu)
 =\operatorname{cf}([I]^{\leq\mu},\subseteq).
\]
In particular, if $|I|\leq\nu$ and $\nu^\mu=\nu$, $J_\mu$ has at most
$\nu$ elements.  More generally a cofinal subposet with at most $\nu$
elements exists exactly when the displayed cofinality is at most $\nu$.
\end{lemma}
\begin{proof}
Starting with a subset $A_0$ of size at most $\mu$, at each successor
stage adjoin one chosen upper bound in $I$ for every subset of the
current set of size less than $\kappa$; take unions at limit stages.
There are at most $\mu^{<\kappa}=\mu$ choices at each stage.  After
$\kappa$ stages the resulting set still has size at most $\mu$ and is
$\kappa$-directed: every set of fewer than $\kappa$ of its elements
appears before one stage, by regularity.  The union of at most $\mu$
members of $J_\mu$ has size at most $\mu$, so the construction gives
an upper bound for it in $J_\mu$.  This proves directedness and also
that $J_\mu$ is cofinal in $[I]^{\leq\mu}$.  Cofinal subposets have
the same cofinality, which proves the formula.  The number of subsets
of size at most $\mu$ in a set of size at most $\nu$ is at most
$\nu^\mu$ (the finite exceptional cases are already excluded).
The last assertion is the definition of cofinality.
\end{proof}

\begin{proposition}[Proved: the comparison map in the cardinal filtration]\label{cat:I3}
Under the hypotheses of Lemma~\ref{cat:I2}, let $D:I\to E$ be a
diagram in a category having the relevant directed colimits.  Set
$X_J=\operatorname{colim}_{i\in J}D(i)$ for $J\in J_\mu$.
The canonical map
\[
 \operatorname{colim}_{J\in J_\mu}X_J
 \longrightarrow\operatorname{colim}_{i\in I}D(i)
\]
is an isomorphism.  Consequently the cardinal bound in the
manuscript's strengthened axiom c) is justified by the two closure
and indexing bounds of Lemma~\ref{cat:I2}.  An equality such as
$\nu^\mu=\nu$ is a sufficient bound, not a necessary condition for an
individual object to admit a small presentation.
\end{proposition}
\begin{proof}
A compatible collection of maps $X_J\to Y$ gives a map $D(i)\to Y$
by choosing any $J$ containing $i$.  Independence follows by choosing
an upper bound of two such $J$'s.  For $i\leq i'$, choose a $J$
containing both; the resulting maps form an $I$-cocone.  Conversely
an $I$-cocone restricts to every $J$ and hence gives the compatible
collection.  These constructions are mutually inverse and natural in
$Y$, proving the isomorphism by the universal property of the colimit.
The bounds are then precisely those in the preceding lemma.
\end{proof}

\begin{proposition}[Refuted: an accessible functor need not preserve all sufficiently large presentability ranks]\label{cat:I4}
The scan-055 assertion
$\exists\kappa_0\ \forall\kappa\geq\kappa_0:
F(E_\kappa)\subseteq F_\kappa$
is false even for an accessible endofunctor of $\mathbf{Set}$.
Take $F(X)=X^{\mathbf N}$.  This functor preserves
$\aleph_1$-filtered colimits, but fails to preserve
$\mu^+$-presentable objects for every infinite cardinal $\mu$ of
countable cofinality.  Such $\mu$ are arbitrarily large.
\end{proposition}
\begin{proof}
A function from $\mathbf N$ into a filtered colimit is represented at
one stage of an $\aleph_1$-filtered diagram, since its countably many
coordinates have a common stage.  Equality of two such functions is
witnessed at one further common stage.  This proves preservation.
In $\mathbf{Set}$ an object is $\lambda$-presentable, for regular
$\lambda$, exactly when its cardinality is less than $\lambda$.
One direction is the same element-and-equality argument; for the other,
write a set as the $\lambda$-directed union of its subsets of size
less than $\lambda$ and apply presentability to its identity.

Now $\mu^{\aleph_0}>\mu$.  For completeness, partition an alleged
list of at most $\mu$ functions $\mathbf N\to\mu$ into countably
many subsets $A_n$, each of cardinality less than $\mu$.  Choose the
$n$th coordinate of a new function outside the $n$th-coordinate
values occurring in $A_n$.  The new function is absent from the list.
Thus a set of size $\mu$ is $\mu^+$-presentable but its image under
$F$ is not.  Taking suprema of increasing countable sequences of
cardinals produces arbitrarily large such $\mu$.
\end{proof}

\begin{proposition}[Proved under stated hypotheses: reconstruction of functors]\label{cat:I5}
If $E$ and $F$ are $\kappa$-accessible and
$T:E\to F$ preserves $\kappa$-filtered colimits and sends
$\kappa$-presentables to $\kappa$-presentables, then the canonical
comparison
$\operatorname{Ind}_\kappa(T|_{E_\kappa})\simeq T$
is a natural isomorphism under Theorem~\ref{cat:I1}'s equivalences.
Conversely every functor obtained by this Ind construction has these
two properties.  The qualification on the rank is essential, as
Proposition~\ref{cat:I4} shows.
\end{proposition}
\begin{proof}
On a presentation $X=\operatorname{colim}_I x_i$ by presentables,
the comparison is the map
$\operatorname{colim}_I T(x_i)\to T(X)$ supplied by the cocone.
It is an isomorphism by preservation.  Using the canonical category
$E_\kappa/X$ makes these maps natural in $X$.  Conversely the Ind
extension is defined by these colimits and preserves them by the
same flattening of categories of elements as in
Theorem~\ref{cat:I1}.  A presentable object is a retract of a
representable, and its image is consequently a retract of a
presentable object of $F$, hence presentable.
\end{proof}

\begin{lemma}[Proved: directed-poset presentations]\label{cat:I7}
Every small $\kappa$-filtered category admits a final functor from a
small $\kappa$-directed poset.  Thus the poset formulation used on
161-2/060 and in Lemma~\ref{cat:I2} loses no presentations.
\end{lemma}
\begin{proof}
Let $D$ be the category.  We construct a poset $P$ with a functor
$u:P\to D$, keeping every principal downset of cardinal less than
$\kappa$.  Start with one isolated vertex for every object of $D$,
with the evident values of $u$.  At a successor stage, for each subset
$S$ of existing vertices of size less than $\kappa$, take the union of
their downsets.  This has fewer than $\kappa$ arrows, by regularity.
Choose a cocone on its image in $D$, and adjoin a new vertex above
that union with this value and these transition maps.  Do not put
relations between vertices added at the same stage.

Also, for each object $d\in D$, pair of old vertices $v,w$, and maps
$a:d\to u(v)$, $b:d\to u(w)$, adjoin an upper vertex whose transition
maps equalize these two maps.  This is possible by taking a cocone on
the union of the two downsets together with a new formal vertex of
value $d$ and arrows of values $a,b$.  The resulting diagram still
has fewer than $\kappa$ arrows.  Its cocone restricts to the old
poset diagram, so it extends $u$ coherently.  Each new principal
downset still has size less than $\kappa$.
Take unions at limit stages and repeat for $\kappa$ stages.
All stages are sets.  Every subset of fewer than $\kappa$ vertices
appears before a common stage and gets an upper bound at the next
one, so $P$ is $\kappa$-directed.  For each $d$, the comma category
$d/u$ is nonempty by the initial vertices, and any two of its objects
are joined by the added equalizing upper vertex.  Hence it is connected.
A cocone on $u$ therefore gives a well-defined map from each $d$,
independent of the choice of an object in $d/u$; the equalizing
vertices verify both independence and compatibility with every arrow
of $D$.  This is exactly finality, directly from the colimit universal
property.
\end{proof}

\begin{theorem}[Proved: the cofinal-rank replacement for the false corollary]\label{cat:I8}
Let $E,F$ be $\kappa$-accessible and let $T:E\to F$ preserve
$\kappa$-filtered colimits.  There are arbitrarily large regular
cardinals $\lambda$ such that $E,F$ are $\lambda$-accessible and
$T$ preserves $\lambda$-filtered colimits and
$\lambda$-presentable objects.
More explicitly, choose small full categories of $\kappa$-presentable
generators $C,D$ for $E,F$, respectively, and for every $c\in C$
choose a $\kappa$-directed presentation of $T(c)$ by objects of $D$.
If $\nu\geq\kappa$ bounds all those indexing sets, any
$\mu\geq\nu$ with $\mu^{<\kappa}=\mu$ and
$\lambda=\mu^+$ has the asserted property.
In particular one may take $\mu=2^\nu$ after increasing $\nu$.
\end{theorem}
\begin{proof}
By Lemma~\ref{cat:I7}, write any object of $E$ as a colimit of a
$\kappa$-directed poset of objects of $C$.
Lemmas~\ref{cat:I2} and~\ref{cat:I3} express it as a
$\lambda$-directed colimit of the colimits of its
$\kappa$-directed subdiagrams of size at most $\mu$.
Each such small colimit is $\lambda$-presentable: its Hom functor
is a limit of at most $\mu<\lambda$ Hom functors of
$\kappa$-presentables, hence of $\lambda$-presentables, and those
limits commute with $\lambda$-filtered colimits of sets.
There is only a set of isomorphism types of such small diagrams.
This proves $\lambda$-accessibility of $E$, and the same proof
applies to $F$.

If $X$ is $\lambda$-presentable, its identity factors through one
stage in this canonical $\lambda$-directed presentation.  Thus it is
a retract of a colimit of a $\kappa$-directed diagram of at most
$\mu$ objects of $C$.  By the bound on the chosen presentations,
every $T(c)$ is $\lambda$-presentable, using the same Hom-limit
argument in $F$.  Applying $T$ to that small directed diagram gives
its colimit, which is again $\lambda$-presentable.  Since a retract
of a presentable object is presentable, $T(X)$ has the required rank.
A $\lambda$-filtered category is $\kappa$-filtered, so preservation
of its colimit already follows from the hypothesis on $T$.
Finally $(2^\nu)^\alpha=2^\nu$ for $0<\alpha<\kappa\leq\nu$,
which verifies the indicated bound and gives arbitrarily large choices.
This replaces ``every sufficiently large rank'' by a proved cofinal
family of ranks.
\end{proof}

\begin{proposition}[Proved: the exact uniform cardinal condition]\label{cat:I9}
For infinite regular cardinals $\kappa<\lambda$, the following are
equivalent:
\begin{enumerate}
\item For every set $X$ of size less than $\lambda$, the poset
$[X]^{<\kappa}$ has a cofinal subset of size less than $\lambda$.
\item In every $\kappa$-directed poset, every subset of size less than
$\lambda$ is contained in a $\kappa$-directed subset of size less than
$\lambda$.
\end{enumerate}
This condition is often written $\kappa\mathrel{\triangleleft}\lambda$.
It is the exact uniform closure condition; a power equality is a
convenient sufficient cardinal bound, not its definition.
The unrestricted closure step used on scan 061 is false for
$\kappa=\aleph_1$ and allowed subdiagram size $\mu=\aleph_\omega$.
\end{proposition}
\begin{proof}
Assume (1) and begin with a subset $A$ of a directed poset of size
less than $\lambda$.  Choose a cofinal family in $[A]^{<\kappa}$
with fewer than $\lambda$ members and adjoin upper bounds only for
members of that family.  Repeat for $\kappa$ stages.  Regularity of
$\lambda>\kappa$ keeps the cardinality less than $\lambda$.
Regularity of $\kappa$ puts every fewer-than-$\kappa$ family of
resulting elements in one stage; the chosen cofinal family at that
stage supplies its upper bound.  This proves (2).
Conversely, apply (2) to the $\kappa$-directed poset
$[X]^{<\kappa}$ and its subset of singletons.
Any $\kappa$-directed subposet containing those singletons is
cofinal, because every member of $[X]^{<\kappa}$ is the union of
fewer than $\kappa$ singletons.  This proves (1).

For the counterexample let $I=[\mu]^{\leq\aleph_0}$, ordered by
inclusion.  It is $\aleph_1$-directed.  If a subposet $J$ of size
at most $\mu$ contains all singletons and is $\aleph_1$-directed,
then it is cofinal in $I$ by the same union argument.
But no family $(A_i)_{i<\mu}$ of countable subsets is cofinal.
Partition its indices into countably many sets $I_n$, each of size
less than $\mu$; this is possible because $\operatorname{cf}\mu=\omega$.
Choose $a_n\in\mu\setminus\bigcup_{i\in I_n}A_i$, since that union
has cardinal less than $\mu$.  The countable set $\{a_n:n\in\mathbf N\}$
is contained in none of the $A_i$, a contradiction.
Thus the missing cardinal hypothesis is real.  Lemma~\ref{cat:I2}
provides the corrected closure estimate, while (1) gives its exact
uniform form.
\end{proof}

\begin{proposition}[Refuted: small presentations need not exhaust presentables]\label{cat:I10}
The question on 161-2/060 whether
$\operatorname{Fil}^{\mu}(E)\ne E_{\mu^+}$ can occur has a positive
answer.  The margin claiming equality from cocompleteness alone is
false.  Let $\kappa=\aleph_1$, $\mu=\aleph_\omega$, let $E$ be the
poset category $\mathcal P(\mu)$ ordered by inclusion, and let
$C=[\mu]^{<\kappa}$ be its full subcategory of countable subsets.
Then $E$ is cocomplete, $E\simeq\operatorname{Ind}_\kappa(C)$, and
the top object $\mu$ is $\mu^+$-presentable but has no
$\kappa$-directed presentation by objects of $C$ with a cofinal
indexing set of size at most $\mu$.
More generally, in $\mathcal P(Z)$ an object $S$ is
$\lambda$-presentable exactly when $|S|<\lambda$, for regular
infinite $\lambda$, and its minimal size of a
$\kappa$-directed cofinal presentation by $<\kappa$-subsets is
\[
 \operatorname{cf}([S]^{<\kappa},\subseteq).
\]
\end{proposition}
\begin{proof}
Colimits in the powerset category are unions.  If $|S|<\lambda$ and
$S$ lies in a $\lambda$-directed union, choose one stage containing
each point of $S$ and then a common upper stage.  Thus
$\operatorname{Hom}(S,-)$ preserves those colimits.
Conversely, $S$ is the $\lambda$-directed union of its subsets of
size less than $\lambda$.  If $|S|\geq\lambda$, none of these
contains $S$, so presentability fails on this presentation.
This proves the presentability assertion.
Every subset of $Z$ is the $\kappa$-directed union of its
$<\kappa$-subsets, so Theorem~\ref{cat:I1} gives the displayed
Ind equivalence.

A $\kappa$-directed family of $<\kappa$-subsets with union $S$
is cofinal in $[S]^{<\kappa}$: for a given $<\kappa$-subset,
choose one member containing each point and take their common upper
bound.  Conversely any cofinal family in that poset is
$\kappa$-directed and has union $S$.  Passing to a cofinal indexing
subset does not change the union.  These observations give exactly
the stated minimal cardinal.
For $S=\mu=\aleph_\omega$ the last paragraph of
Proposition~\ref{cat:I9} proves this cofinality strictly greater
than $\mu$, whereas the first paragraph here proves
$\mu^+$-presentability.  This verifies every claim, including
cocompleteness of the counterexample.
\end{proof}

\begin{theorem}[Proved: the positive cardinal-filtration equality]\label{cat:I11}
Let $E=\operatorname{Ind}_\kappa(C)$, let $\mu\geq\kappa$ be infinite,
and let $\operatorname{Fil}^{\mu}_C(E)$ consist of the objects admitting
a $\kappa$-directed presentation by objects of $C$ with a cofinal index
set of cardinal at most $\mu$.  Always
\[
 \operatorname{Fil}^{\mu}_C(E)\subseteq E_{\mu^+}.
\]
If $\kappa\triangleleft\mu^+$, equivalently
$\operatorname{cf}([\mu]^{<\kappa},\subseteq)\leq\mu$, then equality
holds.  In particular $\mu^{<\kappa}=\mu$ suffices.  This proves exactly
the positive assertion on 161-2/060(b): with its convention
$\kappa=\pi^+$, every $\mu=2^c$ with $c\geq\pi$ satisfies the hypothesis.
At these ranks $E$ is $\mu^+$-accessible and
$E\simeq\operatorname{Ind}_{\mu^+}(E_{\mu^+})$.
If $E$ is cocomplete, the latter accessibility and Ind equivalence hold
at every regular $\lambda\geq\kappa$; this does not force the equality
of the two filtrations at every rank, as Proposition~\ref{cat:I10} shows.
\end{theorem}
\begin{proof}
A cofinal subset of a directed poset is still $\kappa$-directed.  Thus
an object of $\operatorname{Fil}^{\mu}_C(E)$ has a presentation with at
most $\mu$ objects and arrows.  Each object of $C$ is
$\mu^+$-presentable, and a colimit of at most $\mu$ such objects and
arrows remains $\mu^+$-presentable: its Hom functor is the corresponding
limit of at most $\mu$ set-valued functors.  Such limits commute with
$\mu^+$-filtered colimits of sets.  This proves the inclusion without
a cardinal-arithmetic hypothesis.

Under the sharp-cardinal hypothesis, Theorem~\ref{cat:I9}'s closure
property permits the construction of $J_\mu$ and the decomposition in
Lemmas~\ref{cat:I2}--\ref{cat:I3} with precisely the same bound $\mu$.
Consequently every $X\in E_{\mu^+}$ is a retract of some
$Y=\operatorname{colim}_{j\in J}c_j$, where $J$ is $\kappa$-directed
and $|J|\leq\mu$.  Indeed $X$ is a $\mu^+$-directed colimit of these
objects, so its identity factors through one of them.
It remains to remove this retract; it must not simply be omitted.

Write $i:X\to Y$, $p:Y\to X$, $pi=1$, $e=ip$, and let
$a_j:c_j\to Y$ be the presentation maps.  By
$\kappa$-presentability, for each $j$ there are $k(j)\in J$ and
$f_j:c_j\to c_{k(j)}$ satisfying $a_{k(j)}f_j=ea_j$.
Put $q_j=pa_j$.  Inside $C/X$ take the subcategory $D$ generated by
the presentation arrows and these additional arrows $f_j$; the latter
belong to $C/X$ because $q_{k(j)}f_j=q_j$.  There are at most $\mu$
objects and arrows after closing under finite composition.  Its
tautological diagram has colimit $X$: a cocone on the original
presentation is a map $g:Y\to Z$, and the additional conditions are
exactly $ge=g$.  Such maps are uniquely $hp$ for $h:X\to Z$.

We enlarge $D$ to a $\kappa$-filtered subcategory $D'\subset C/X$
with at most $\mu$ arrows, preserving this colimit.  At each of
$\kappa$ stages choose a cofinal family, of cardinal at most $\mu$,
of the subsets of the current arrow set having size less than
$\kappa$.  For each chosen subset add a cocone in $C/X$ on its
underlying graph.  Such a cocone exists: first dominate its fewer than
$\kappa$ objects, then equalize its fewer than $\kappa$ arrow
conditions, using $\kappa$-filteredness of $C/X$.
For every new object $(c,q)$ also add one arrow from it to an original
object $(c_j,q_j)$.  This is possible because $iq:c\to Y$ factors
through some $a_j$, and applying $p$ gives the required factorization
of $q$.  Close under composition and take unions at limit stages.
All stages and their union have at most $\mu$ arrows.
Regularity of $\kappa$ ensures that every diagram with fewer than
$\kappa$ arrows in the union occurs at one stage and acquires a cocone
at the next.  Hence $D'$ is $\kappa$-filtered.
Its colimit remains $X$: restriction of any cocone to $D$ gives its
unique map $X\to Z$, and the chosen outgoing arrows to $D$ force
the component at every added object to be the corresponding composite.

Finally the directed-poset replacement of Lemma~\ref{cat:I7} can be
made with at most $\mu$ vertices.  In that construction choose a
cofinal family of fewer-than-$\kappa$ subsets at each stage instead
of all such subsets; there are at most $\mu$ choices by the same
sharp-cardinal hypothesis.  The additional pairs used for comma-category
connectivity are also at most $\mu$, and there are $\kappa\leq\mu$
stages.  This yields a $\kappa$-directed presentation of $X$ of size
at most $\mu$, proving equality.

There is only a set of diagrams of size at most $\mu$ in the small
category $C$.  The decomposition by $J_\mu$ therefore proves
$\mu^+$-accessibility, and Theorem~\ref{cat:I1} gives the displayed
Ind equivalence.  If $\mu=2^c$ and $c\geq\pi$, then for every
$\alpha<\kappa=\pi^+$ one has
$(2^c)^\alpha=2^{c\alpha}=2^c$ when $\alpha>0$; hence
$\mu^{<\kappa}=\mu$.
For the final assertion, a small dense category of
$\kappa$-presentables is still dense and consists of
$\lambda$-presentables.  The cocomplete clause of
Theorem~\ref{cat:I1} applies at every such regular $\lambda$.
\end{proof}


\begin{theorem}[Proved: the functor equivalences with their exact meaning]\label{cat:I12}
For accessible categories $E,F$, the following are equivalent.
\begin{enumerate}
\item $T:E\to F$ is accessible.
\item At some common accessibility rank $\lambda$, it is the Ind
extension of a functor $E_\lambda\to F_\lambda$.
\item It is the pointwise left Kan extension of its restriction to a
small full subcategory of $E$.
\end{enumerate}
In (3), ``pointwise'' means that the canonical colimits over $C/X$
exist and compute the extension; this specifies the meaning of the
manuscript's $i_!^F$ on scan 055.
Moreover, if $E$ is $\kappa$-accessible and a locally small category
$F$ has $\kappa$-filtered colimits, restriction is an equivalence
\[
 \operatorname{Fun}_{\kappa\text{-filt}}(E,F)
   \xrightarrow{\sim}\operatorname{Fun}(E_\kappa,F).
\]
Here the left side consists of functors preserving those colimits,
with all natural transformations.  No presentability condition on
the values in $F$ is needed for this last equivalence.
\end{theorem}
\begin{proof}
Choose a small skeleton of $E_\kappa$.  The extension of a functor
$h:E_\kappa\to F$ is
\[
 \overline h(X)=\operatorname{colim}_{(c\to X)\in E_\kappa/X}h(c).
\]
The indexing category is $\kappa$-filtered by
Theorem~\ref{cat:I1}, so the colimit exists.  Its universal property
identifies maps $\overline h(X)\to Z$ with natural transformations
\[
 \operatorname{Hom}_E(-,X)|_{E_\kappa}
       \longrightarrow \operatorname{Hom}_F(h(-),Z).
\]
For a $\kappa$-filtered colimit of $X$'s the presheaf on the left is
their pointwise colimit.  Mapping this colimit to the right-hand
presheaf gives the limit of the corresponding sets of natural
transformations.  Thus $\overline h$ preserves the required colimits.
At a presentable object $c$ the slice has terminal object $1_c$, so
the restriction is $h$.  Conversely any preserving functor is
computed by these canonical presentations.  The same colimit
property extends every natural transformation uniquely.  This proves
the displayed equivalence and the pointwise Kan assertion at this rank.

For an accessible $T$, first choose a cardinal $\nu$ above given
accessibility ranks for $E,F$ and a preservation rank for $T$.
Theorem~\ref{cat:I11}, applied separately to $E$ and $F$, makes
$(2^\nu)^+$ a common accessibility rank, and $T$ preserves filtered
colimits at that rank because it is larger than its preservation rank.
Theorem~\ref{cat:I8} now supplies a common rank
at which $T$ also preserves presentables; Proposition~\ref{cat:I5}
then gives (2).  Condition (2) implies (1) and (3) by the preceding
paragraphs.
Suppose (3) holds for a small full category $C$.  Each object of an
accessible category is presentable at some rank: write it as a small
filtered colimit of presentables and take a regular cardinal larger
than the indexing category's arrow set.  The Hom-limit calculation
used in Theorem~\ref{cat:I11} proves this assertion directly.
Since $C$ is small, choose one regular $\lambda$ making all its
objects $\lambda$-presentable and such that $E,F$ admit
$\lambda$-filtered colimits.  The same natural-transformation formula,
now with $C$ in place of $E_\kappa$, shows that the pointwise Kan
extension preserves those colimits.  This is (1).
\end{proof}


\subsection{The relative classifying topos}
All topoi in the following statement are Grothendieck topoi over the
fixed universe.  We use a consistent inverse-image convention for
2-morphisms.  Let $\mathbf B_I$ denote the classifier of a small
$I$-indexed family of objects.  For one sort it is
$[\mathbf{FinSet},\mathbf{Set}]$; for several sorts it is the
classifying topos of the empty theory with those sorts.  The notation
$\widehat{(\mathbf{FinSet})^{\mathrm{op}}}^{\,I}$ on scan 112 must
mean a product in the 2-category of topoi if it denotes this
classifier; the ordinary product of the underlying categories is
not the classifier of several independent objects.

\begin{theorem}[Proved, using the opened classifying-topos existence theorem]\label{cat:I6}
Let $T$ be a small geometric theory with classifying topos $B$, and
let $(b_i)_{i\in I}$ be specified sorts of its universal model.
For a Grothendieck topos $X$ and objects $E_i\in X$ there is a
Grothendieck topos $X[T;E]\to X$ such that, for every
$g:Y\to X$, geometric morphisms $Y\to X[T;E]$ over $X$ classify
$T$-models $M$ in $Y$ equipped with specified isomorphisms
$M_i\simeq g^*E_i$.
With $b:B\to\mathbf B_I$ and $e:X\to\mathbf B_I$ the morphisms
classifying the two families, we have a 2-pullback
\[
 X[T;E]\simeq B\times_{\mathbf B_I}X.
\]
This includes small finite-limit theories, as requested on
cote 161-2/112.
\end{theorem}
\begin{proof}
Choose a small generating full subcategory $C$ of $X$, enlarge it to
contain every $E_i$, and close it under finite limits.  This remains
small.  Equip $C$ with the induced canonical topology $J$; the
comparison functor gives $X\simeq\operatorname{Sh}(C,J)$.
The geometric theory $T_C$ has a sort for every object of $C$, a
function for every arrow, the equations for identities and composition,
the finite-limit axioms, and the following axiom for each covering
family $(c_a\to c)_a$:
\[
 x=x\ \vdash_{x:c}\ \bigvee_a\exists y:c_a\ (f_a(y)=x).
\]
Its models in $Y$ are exactly the left exact $J$-continuous functors
$C\to Y$, equivalently geometric morphisms $Y\to X$.
This identification, including morphisms, is the site-classification
result in Makkai--Reyes~\cite[Propositions 9.1.4--9.1.5]{catMakkaiReyes1977}.

Form the disjoint union of the signatures and axioms of $T_C$ and
$T$.  For every $i$, add two function symbols
$u_i:M_i\to F(E_i)$ and $v_i:F(E_i)\to M_i$, with equations
$v_i u_i=1$ and $u_i v_i=1$.  These are geometric axioms, so the
resulting small geometric theory $T'$ has a classifying topos by
\cite[Theorem 9.1.1]{catMakkaiReyes1977}.  Define
$X[T;E]=\operatorname{Set}[T']$.  Its universal $T_C$-model gives
its morphism to $X$.

By construction, a model of $T'$ in $Y$ is exactly a triple
$(g,M,(u_i)_i)$ as in the statement.  Morphisms of models are
exactly the compatible pairs of morphisms with the indicated
identifications, because the inverse equations impose no extra choices.
The equivalence is natural under inverse-image functors: those preserve
finite limits, arbitrary unions, and images, hence all the displayed
axioms.  The hom-category of triples is the pseudopullback of
$\operatorname{Hom}(Y,B)$ and $\operatorname{Hom}(Y,X)$ over
$\operatorname{Hom}(Y,\mathbf B_I)$.  This is precisely the
2-pullback universal property, and fixing $g$ gives the relative
universal property.  A finite-limit theory is geometric by expressing
its specified limits as existence-and-uniqueness axioms, so it is
included.
\end{proof}
